\label{SI:priorcov}
{\color{blue}
To assess the sensitivity of the EnKF to the choice of prior parameter covariance width, the nominal covariance matrix $\Sigma_0^\theta$ was scaled by factors of $0.5\times$, $1.0\times$, $1.5\times$, and $2.0\times$, and the filter was re-run for all six datasets using the best-tuned ensemble size for each. Figures~\ref{fig:priorcov_A_flask}--\ref{fig:priorcov_B_pl40} show the resulting state trajectories alongside experimental data for each dataset. The $1.0\times$ case corresponds to the nominal results presented in the main body and is stable by definition. Table~\ref{tab:priorcov_stability} summarises per-state stability: a state is flagged as unstable (\ding{55}) if its trajectory deviates from the $1.0\times$ reference by more than 50\% of the experimental maximum for that state at any point in the simulation. The $0.5\times$ setting remains stable across all datasets. Instabilities arise at $2.0\times$ for several metabolite states, particularly in the Cell Line B datasets, while the $1.5\times$ setting only affects asparagine in the Cell Line A Bioreactor 32\textdegree C condition.
}

\begin{figure}[H]
    \centering
    \includegraphics[width=0.95\textwidth]{Figs/prior_width_profiles_CHO_T127_flask_PMJ.png}
    \caption{\rev{Effect of prior covariance width on EnKF state estimation for Cell Line A Shake Flask. Trajectories are shown for $0.5\times$ (dark solid), $1.0\times$ (solid), $1.5\times$ (dashed), and $2.0\times$ (dotted) nominal covariance, alongside experimental measurements (markers, $\pm$standard deviation).}}
    \label{fig:priorcov_A_flask}
\end{figure}

\begin{figure}[H]
    \centering
    \includegraphics[width=0.95\textwidth]{Figs/prior_width_profiles_CHO_T127_SNS_36.5.png}
    \caption{\rev{Effect of prior covariance width on EnKF state estimation for Cell Line A Bioreactor 36.5\textdegree C. Trajectories are shown for $0.5\times$ (dark solid), $1.0\times$ (solid), $1.5\times$ (dashed), and $2.0\times$ (dotted) nominal covariance, alongside experimental measurements (markers, $\pm$standard deviation).}}
    \label{fig:priorcov_A_36}
\end{figure}

\begin{figure}[H]
    \centering
    \includegraphics[width=0.95\textwidth]{Figs/prior_width_profiles_CHO_T127_SNS_32.png}
    \caption{\rev{Effect of prior covariance width on EnKF state estimation for Cell Line A Bioreactor 32\textdegree C. Trajectories are shown for $0.5\times$ (dark solid), $1.0\times$ (solid), $1.5\times$ (dashed), and $2.0\times$ (dotted) nominal covariance, alongside experimental measurements (markers, $\pm$standard deviation).}}
    \label{fig:priorcov_A_32}
\end{figure}

\begin{figure}[H]
    \centering
    \includegraphics[width=0.95\textwidth]{Figs/prior_width_profiles_CHO_GS46_F_C_Inv.png}
    \caption{\rev{Effect of prior covariance width on EnKF state estimation for Cell Line B Feed C. Trajectories are shown for $0.5\times$ (dark solid), $1.0\times$ (solid), $1.5\times$ (dashed), and $2.0\times$ (dotted) nominal covariance, alongside experimental measurements (markers, $\pm$standard deviation).}}
    \label{fig:priorcov_B_C}
\end{figure}

\begin{figure}[H]
    \centering
    \includegraphics[width=0.95\textwidth]{Figs/prior_width_profiles_CHO_GS46_F_all.png}
    \caption{\rev{Effect of prior covariance width on EnKF state estimation for Cell Line B Feed U. Trajectories are shown for $0.5\times$ (dark solid), $1.0\times$ (solid), $1.5\times$ (dashed), and $2.0\times$ (dotted) nominal covariance, alongside experimental measurements (markers, $\pm$standard deviation).}}
    \label{fig:priorcov_B_U}
\end{figure}

\begin{figure}[H]
    \centering
    \includegraphics[width=0.95\textwidth]{Figs/prior_width_profiles_CHO_GS46_F_all_pl40.png}
    \caption{\rev{Effect of prior covariance width on EnKF state estimation for Cell Line B Feed U plus 40\%. Trajectories are shown for $0.5\times$ (dark solid), $1.0\times$ (solid), $1.5\times$ (dashed), and $2.0\times$ (dotted) nominal covariance, alongside experimental measurements (markers, $\pm$standard deviation).}}
    \label{fig:priorcov_B_pl40}
\end{figure}

% cell-colour shorthands (local to this file)
\definecolor{stabgreen}{RGB}{163,210,163}   % muted sage green — stable
\definecolor{unstabred}{RGB}{218,112,112}   % muted crimson — unstable
\newcommand{\G}{\cellcolor{stabgreen}\ding{51}}   % tick
\newcommand{\R}{\cellcolor{unstabred}\ding{55}}   % cross

\begin{table}[ht]
\centering
\renewcommand{\arraystretch}{1.3}
\small
\setlength{\tabcolsep}{3pt}
\caption{\rev{Per-state EnKF stability under different prior covariance widths for all six datasets. Green cells indicate a stable trajectory; red cells indicate that the trajectory deviates from the $1.0\times$ baseline by more than 50\% of the experimental maximum for that state at any simulation point. The $1.0\times$ column is the nominal case used throughout the paper and is stable by definition.}}
\label{tab:priorcov_stability}
\arrayrulecolor{black}
\adjustbox{max width=\textwidth}{%
\begin{tabular}{|l|p{1.6cm}|>{\centering\arraybackslash}p{2.4cm}|>{\centering\arraybackslash}p{2.4cm}|>{\centering\arraybackslash}p{2.4cm}|>{\centering\arraybackslash}p{2.4cm}|}
\hline
\rowcolor{gray!30}
\textbf{Dataset} & \textbf{State} &
\textbf{\makecell{0.5 $\times$ Baseline\\Prior Width}} &
\textbf{\makecell{1.0 $\times$ Baseline\\Prior Width}} &
\textbf{\makecell{1.5 $\times$ Baseline\\Prior Width}} &
\textbf{\makecell{2.0 $\times$ Baseline\\Prior Width}} \\
\hline
%──────────────────────────────────────────────────────
% Cell Line A — Shake Flask  (all stable)
%──────────────────────────────────────────────────────
\multirow{8}{*}{\makecell[l]{Cell Line A\\Shake Flask}}
 & Xv  & \G & \G & \G & \G \\ \hhline{|~|-----|}
 & mAb & \G & \G & \G & \G \\ \hhline{|~|-----|}
 & Glc & \G & \G & \G & \G \\ \hhline{|~|-----|}
 & Amm & \G & \G & \G & \G \\ \hhline{|~|-----|}
 & Gln & \G & \G & \G & \G \\ \hhline{|~|-----|}
 & Lac & \G & \G & \G & \G \\ \hhline{|~|-----|}
 & Glu & \G & \G & \G & \G \\ \hhline{|~|-----|}
 & Asn & \G & \G & \G & \G \\
\hline
%──────────────────────────────────────────────────────
% Cell Line A — Bioreactor 36.5 °C  (2.0× Glc unstable)
%──────────────────────────────────────────────────────
\multirow{8}{*}{\makecell[l]{Cell Line A\\Bioreactor 36.5\textdegree C}}
 & Xv  & \G & \G & \G & \G \\ \hhline{|~|-----|}
 & mAb & \G & \G & \G & \G \\ \hhline{|~|-----|}
 & Glc & \G & \G & \G & \R \\ \hhline{|~|-----|}
 & Amm & \G & \G & \G & \G \\ \hhline{|~|-----|}
 & Gln & \G & \G & \G & \G \\ \hhline{|~|-----|}
 & Lac & \G & \G & \G & \G \\ \hhline{|~|-----|}
 & Glu & \G & \G & \G & \G \\ \hhline{|~|-----|}
 & Asn & \G & \G & \G & \G \\
\hline
%──────────────────────────────────────────────────────
% Cell Line A — Bioreactor 32 °C  (1.5× Asn; 2.0× Lac+Asn unstable)
%──────────────────────────────────────────────────────
\multirow{8}{*}{\makecell[l]{Cell Line A\\Bioreactor 32\textdegree C}}
 & Xv  & \G & \G & \G & \G \\ \hhline{|~|-----|}
 & mAb & \G & \G & \G & \G \\ \hhline{|~|-----|}
 & Glc & \G & \G & \G & \G \\ \hhline{|~|-----|}
 & Amm & \G & \G & \G & \G \\ \hhline{|~|-----|}
 & Gln & \G & \G & \G & \G \\ \hhline{|~|-----|}
 & Lac & \G & \G & \G & \R \\ \hhline{|~|-----|}
 & Glu & \G & \G & \G & \G \\ \hhline{|~|-----|}
 & Asn & \G & \G & \R & \R \\
\hline
%──────────────────────────────────────────────────────
% Cell Line B — Feed C  (2.0× Glc+Amm+Gln+Glu+Asn unstable)
%──────────────────────────────────────────────────────
\multirow{8}{*}{\makecell[l]{Cell Line B\\Feed C}}
 & Xv  & \G & \G & \G & \G \\ \hhline{|~|-----|}
 & mAb & \G & \G & \G & \G \\ \hhline{|~|-----|}
 & Glc & \G & \G & \G & \R \\ \hhline{|~|-----|}
 & Amm & \G & \G & \G & \R \\ \hhline{|~|-----|}
 & Gln & \G & \G & \G & \R \\ \hhline{|~|-----|}
 & Lac & \G & \G & \G & \G \\ \hhline{|~|-----|}
 & Glu & \G & \G & \G & \R \\ \hhline{|~|-----|}
 & Asn & \G & \G & \G & \R \\
\hline
%──────────────────────────────────────────────────────
% Cell Line B — Feed U  (2.0× Amm+Gln+Lac+Asn unstable)
%──────────────────────────────────────────────────────
\multirow{8}{*}{\makecell[l]{Cell Line B\\Feed U}}
 & Xv  & \G & \G & \G & \G \\ \hhline{|~|-----|}
 & mAb & \G & \G & \G & \G \\ \hhline{|~|-----|}
 & Glc & \G & \G & \G & \G \\ \hhline{|~|-----|}
 & Amm & \G & \G & \G & \R \\ \hhline{|~|-----|}
 & Gln & \G & \G & \G & \R \\ \hhline{|~|-----|}
 & Lac & \G & \G & \G & \R \\ \hhline{|~|-----|}
 & Glu & \G & \G & \G & \G \\ \hhline{|~|-----|}
 & Asn & \G & \G & \G & \R \\
\hline
%──────────────────────────────────────────────────────
% Cell Line B — Feed U plus 40%  (2.0× Amm+Gln+Asn unstable)
%──────────────────────────────────────────────────────
\multirow{8}{*}{\makecell[l]{Cell Line B\\Feed U plus 40\%}}
 & Xv  & \G & \G & \G & \G \\ \hhline{|~|-----|}
 & mAb & \G & \G & \G & \G \\ \hhline{|~|-----|}
 & Glc & \G & \G & \G & \G \\ \hhline{|~|-----|}
 & Amm & \G & \G & \G & \R \\ \hhline{|~|-----|}
 & Gln & \G & \G & \G & \R \\ \hhline{|~|-----|}
 & Lac & \G & \G & \G & \G \\ \hhline{|~|-----|}
 & Glu & \G & \G & \G & \G \\ \hhline{|~|-----|}
 & Asn & \G & \G & \G & \R \\
\hline
\end{tabular}}
\end{table}
